\chapter{Kết luận và Hướng phát triển}
\label{Chapter5}

Chương này tổng kết những kết quả chính của đề tài, đánh giá phạm vi mà Yoca đã giải quyết và trình bày các hướng phát triển có thể mở rộng giá trị của hệ thống. Các kết quả được xem xét trên cả hai phương diện: trải nghiệm phân tích dành cho người dùng và nền tảng kỹ thuật phục vụ việc thu thập, tổ chức và diễn giải dữ liệu Solana.

\section{Tổng kết các kết quả đạt được}

Yoca đã hình thành một hành trình phân tích liên tục từ thị trường đến dữ liệu ví. Người dùng có thể bắt đầu tại trang tổng quan để nhận biết token hoặc pool đáng chú ý, mở trang chi tiết để xem giá, thanh khoản, holder và giao dịch, sau đó tiếp tục đến các ví liên quan để quan sát danh mục tài sản và hành vi theo thời gian. Thanh tìm kiếm và các liên kết giữa những đối tượng giúp quá trình này diễn ra trong cùng một hệ thống, giữ được ngữ cảnh khi người dùng chuyển từ góc nhìn tổng quan sang dữ liệu chi tiết.

Phân tích ví là một trong những kết quả trọng tâm của đề tài. Yoca tổng hợp giá trị danh mục, lịch sử số dư, khối lượng giao dịch, PnL và các vị thế token theo nhiều khoảng thời gian. Lịch sử swap và transfer được chuẩn hóa để người dùng theo dõi hoạt động của ví và mở chi tiết từng giao dịch khi cần đối chiếu. Chức năng so sánh nhiều ví bổ sung góc nhìn tương quan giữa hiệu quả, mức độ hoạt động và cơ cấu tài sản, qua đó hỗ trợ việc nhận biết những khác biệt đáng chú ý giữa các địa chỉ.

Mô-đun Wash Trading mở rộng phạm vi phân tích từ từng giao dịch sang cấu trúc quan hệ giữa nhiều ví. Dữ liệu chuyển token được biểu diễn bằng đồ thị có hướng; mỗi cạnh thể hiện luồng chuyển từ ví gửi đến ví nhận và có thể được gộp theo cùng chiều giữa một cặp ví để giảm chồng lấn khi hiển thị. Các đặc trưng về chu trình, khoảng thời gian trung bình giữa các giao dịch liên tiếp, độ tương đồng lượng token, self-loop, hubness và khối lượng tương đối được kết hợp để tạo điểm nghi vấn cho ví và điểm rủi ro tổng thể của token. Kết quả được trình bày cùng đồ thị, danh sách ví nghi vấn và các chu trình nghi vấn, giúp người dùng quan sát cơ sở hình thành tín hiệu trước khi tiếp tục đối chiếu dữ liệu giao dịch.

Trợ lý AI được tích hợp vào các miền token, ví và Wash Trading để diễn giải dữ liệu theo ngữ cảnh. Máy chủ chuẩn bị các chỉ số và dữ liệu liên quan trước khi gửi yêu cầu đến mô hình; các tham số đầu vào được kiểm tra bằng schema, còn phản hồi từ mô hình được yêu cầu ở dạng JSON để đọc và chuẩn hóa các trường cần thiết. Riêng với Wash Trading, ngữ cảnh gồm mint, khung thời gian, cấu hình chấm điểm và kết quả phân tích hiện tại; khi mô hình không khả dụng hoặc phản hồi không hợp lệ, hệ thống sử dụng phần phân tích cục bộ làm phương án thay thế. Các chức năng watchlist, nhãn ví và cảnh báo hỗ trợ người dùng lưu lại đối tượng quan tâm, theo dõi sự kiện và duy trì quá trình phân tích qua nhiều lần sử dụng. Giao diện tiếng Việt cùng cách trình bày số, ngày giờ và giá trị quy đổi giúp hệ thống gần gũi hơn với người dùng trong nước.

Ở lớp kỹ thuật, Yoca xây dựng một mô hình tích hợp dữ liệu theo miền. CoinGecko, GeckoTerminal và Birdeye cung cấp dữ liệu thị trường; Helius phục vụ danh mục ví, giao dịch và sự kiện on-chain; Mobula cung cấp các chỉ số phân tích ví và lịch sử tổng tài sản; Zerion đảm nhiệm lịch sử ở cấp token; Moralis bổ sung metadata cho một số luồng dữ liệu. Các nguồn này được đặt sau lớp kiểm tra và chuẩn hóa chung, nhờ đó phần giao diện không phụ thuộc trực tiếp vào cấu trúc phản hồi riêng của từng dịch vụ.

Cơ sở dữ liệu vừa lưu thông tin nghiệp vụ, vừa tái sử dụng dữ liệu đã đồng bộ. Dữ liệu dạng snapshot được quản lý theo thời gian hiệu lực, còn lịch sử giao dịch và số dư được theo dõi theo khoảng đã tải. Cách tổ chức này giảm số lần truy vấn lặp, hỗ trợ phân trang dữ liệu dài và cho phép từng nhóm thông tin được cập nhật theo nhịp riêng. Hợp đồng kiểu giữa client và server, kiểm tra dữ liệu ở runtime và các ràng buộc trong PostgreSQL tạo thành nhiều lớp bảo vệ cho quá trình trao đổi và lưu trữ dữ liệu.

Đề tài cũng hoàn thiện quy trình kiểm tra và triển khai sản phẩm. Các bộ kiểm thử phía client và server bao phủ những miền quan trọng như xác thực, thanh toán, cảnh báo, dữ liệu ví, provider adapter và AI. Lần chạy được ghi nhận trong báo cáo đạt 178/178 test phía client và 204/204 test phía server. GitHub Actions thực hiện typecheck, lint, test và production build trước khi kích hoạt quá trình triển khai. Giao diện và máy chủ được vận hành thành hai dịch vụ độc lập trên Render, kết nối với PostgreSQL trên Supabase và phục vụ tại tên miền của Yoca.

\section{Giới hạn của đề tài}

Khả năng phân tích của Yoca chịu ảnh hưởng bởi độ phủ dữ liệu từ các dịch vụ bên ngoài. Token mới, tài sản có thanh khoản thấp hoặc giao dịch có cấu trúc phức tạp có thể thiếu metadata, giá lịch sử hoặc thông tin dẫn xuất cần thiết. Quota và phạm vi lịch sử của từng nhà cung cấp cũng tác động đến tần suất cập nhật. Hệ thống sử dụng cache, kiểm tra dữ liệu và phân công nguồn theo từng miền để duy trì kết quả nhất quán trong phạm vi dữ liệu nhận được.

Các chỉ số PnL, nội dung AI và điểm Wash Trading được sử dụng như công cụ hỗ trợ phân tích. Giá trị của chúng phụ thuộc vào khoảng thời gian, dữ liệu giá và tập giao dịch được quan sát. Đối với Wash Trading, các đặc trưng và trọng số chỉ làm nổi bật các cấu trúc có dấu hiệu bất thường trên đồ thị; người dùng cần đối chiếu thêm giao dịch gốc, bối cảnh thị trường và nguồn dữ liệu liên quan trước khi đưa ra nhận định cuối cùng. Kết quả từ nguồn demo-fallback chỉ minh họa luồng giao diện và không được xem là kết luận từ dữ liệu on-chain thật. Hoạt động đánh giá của đề tài tập trung vào tính đúng đắn của luồng xử lý và các chức năng cốt lõi.

Ngoài kết quả về chức năng, đề tài đã giải quyết các vấn đề quan trọng khi làm việc với dữ liệu Web3. Cơ chế điều phối yêu cầu, luân phiên khóa truy cập và thử lại có kiểm soát giúp hạn chế tác động của giới hạn truy vấn. Chiến lược lưu đệm theo phạm vi thời gian cho phép tái sử dụng dữ liệu đã có và chỉ truy xuất phần còn thiếu. Đối với giá trị danh mục lịch sử, hệ thống tái dựng số dư tại từng thời điểm và kết hợp với dữ liệu giá tương ứng, qua đó tạo cơ sở cho biểu đồ số dư và PnL.

Hướng phát triển đầu tiên là mở rộng chiều sâu của dữ liệu và khả năng theo dõi chất lượng nguồn. Khi hệ thống phục vụ nhiều người dùng hơn, các chỉ số về thời gian phản hồi, tỷ lệ lỗi, mức sử dụng quota và hiệu quả cache có thể được tổng hợp để điều chỉnh chiến lược cập nhật theo từng provider. Dữ liệu lịch sử cũng có thể được lưu trữ theo các lớp thời gian phù hợp với nhu cầu truy vấn, giúp các phân tích dài hạn duy trì tốc độ phản hồi ổn định.

Đối với Wash Trading, hệ thống có thể bổ sung tập dữ liệu được gán nhãn và mở rộng thực nghiệm trên nhiều token, pool cùng giai đoạn thị trường khác nhau. Các đặc trưng hiện có tạo nền tảng cho việc đánh giá thêm những yếu tố như biến động giá, thay đổi thanh khoản, tần suất lặp lại giữa các cụm ví và hành vi qua nhiều pool. Khi dữ liệu đánh giá đủ lớn, các cấu hình GCN, GAT và GraphSAGE có thể được phát triển từ bộ trọng số heuristic hiện tại thành mô hình học trên đồ thị hoặc được hiệu chỉnh bằng phương pháp thống kê để tăng độ tin cậy của kết quả.

Các chức năng AI có thể tăng khả năng giải thích bằng cách liên kết từng nhận định với giao dịch, vị trí trên biểu đồ hoặc chỉ số đã sử dụng. Điều này giúp người dùng chuyển trực tiếp từ phần diễn giải đến dữ liệu cần kiểm tra. Cùng với đó, Yoca có thể mở rộng không gian cá nhân thành một quy trình nghiên cứu hoàn chỉnh hơn, nơi watchlist, nhãn ví, cảnh báo và lịch sử phân tích được kết nối theo từng chủ đề mà người dùng quan tâm.

Nhìn chung, đề tài đã xây dựng được một nền tảng phân tích dữ liệu Solana có luồng sử dụng rõ ràng, kết hợp dữ liệu thị trường với phân tích ví, giao dịch và quan hệ đồ thị. Quá trình thực hiện giúp nhóm chúng em giải quyết đồng thời các bài toán tích hợp nhiều nguồn, quản lý dữ liệu lịch sử, trực quan hóa thông tin và hỗ trợ diễn giải kết quả. Những thành phần đã xây dựng tạo cơ sở để Yoca tiếp tục phát triển thành một công cụ phân tích on-chain có chiều sâu và phù hợp hơn với nhu cầu của người dùng Việt Nam.
